\documentclass[10pt, twocolumn]{article}

% ── Packages ─────────────────────────────────────────────────────────────────
\usepackage[T1]{fontenc}
\usepackage[utf8]{inputenc}
\usepackage[margin=1in]{geometry}
\usepackage{amsmath, amssymb}
\usepackage{graphicx}
\usepackage{booktabs}
\usepackage{multirow}
\usepackage{array}
\usepackage{url}
\usepackage{hyperref}
\usepackage{cite}
\usepackage{balance}
\usepackage{microtype}
\usepackage{xcolor}
\usepackage{float}

% ── Title ─────────────────────────────────────────────────────────────────────
\title{\textbf{Hybrid Teacher Ensemble Knowledge Distillation with\\
Ordinal-Aware Loss for Low-Data Diabetic Retinopathy Grading}}

\author{
  Author 1\textsuperscript{a} \quad
  Author 2\textsuperscript{a} \quad
  Author 3\textsuperscript{b} \quad
  Author 4\textsuperscript{b,*}\\[4pt]
  \textsuperscript{a}\small{[Department], [University], [City], [Country]}\\
  \textsuperscript{b}\small{[Department], [University], [City], [Country]}\\
  \textsuperscript{*}\small{Corresponding author. E-mail: [corresponding@email.com]}
}

\date{}

\begin{document}
\maketitle

% ── Abstract ──────────────────────────────────────────────────────────────────
\begin{abstract}
Automated grading of diabetic retinopathy (DR) from fundus photographs is a
clinically critical task constrained by two persistent challenges: the scarcity of
expert-annotated training data and the ordinal structure of the five-class severity
scale, which standard cross-entropy losses fail to encode. This paper proposes
\textit{TS\_ConvNeXtTiny\_Residual}, a teacher--student knowledge distillation
framework designed specifically for five-class DR severity grading under low-data
conditions. The framework employs an ensemble of two structurally diverse CNN-ViT
hybrid teacher networks---a custom Baseline Hybrid with confidence-gated fusion and
an Advanced ResNet50+ViT-B/16 model with spectral normalization---whose softmax
predictions are combined via equal-weight averaging to form a well-calibrated
composite soft-label signal. A compact ConvNeXt~Tiny student network is trained
using a four-component objective integrating: (i)~class-weighted cross-entropy for
imbalance handling, (ii)~multi-temperature KL-divergence distillation at
$T \in \{2.0, 4.0\}$ for simultaneous transfer of fine-grained and coarse semantic
knowledge, (iii)~an ordinal cumulative distribution function (CDF) loss that aligns
the training objective with the quadratic weighted kappa (QWK) evaluation metric,
and (iv)~a non-inferiority loss with a linearly scheduled margin that prevents the
student from regressing below teacher-ensemble performance on individual samples.
Evaluated on the APTOS~2019 benchmark across three random seeds, the proposed
system achieves a mean test QWK of $\mathbf{0.9186 \pm 0.0010}$ and a mean
accuracy of $\mathbf{85.47\% \pm 0.31\%}$, surpassing all single-teacher baselines
and alternative teacher--student configurations examined in this study.
\end{abstract}

\textbf{Keywords:} diabetic retinopathy grading; knowledge distillation; CNN-ViT hybrid;
ordinal CDF loss; teacher ensemble; low-data learning; APTOS 2019

% ── 1. Introduction ───────────────────────────────────────────────────────────
\section{Introduction}

Diabetic retinopathy (DR) is the most prevalent microvascular complication of
diabetes mellitus and one of the leading causes of preventable vision loss and
blindness worldwide~\cite{yau2012}. The International Diabetes Federation reports
that approximately 537 million adults were living with diabetes in 2021, a figure
projected to rise to 783 million by 2045~\cite{idf2021}. Among diabetic
individuals, DR affects approximately one-third, with a significant proportion
progressing to sight-threatening stages if not detected in a timely manner.

Conventional DR screening relies on trained ophthalmologists who manually examine
retinal fundus images and assign a severity grade according to standardized clinical
scales. The International Clinical Diabetic Retinopathy (ICDR) severity scale,
which classifies retinopathy into five ordinal grades from no apparent retinopathy
(Grade~0) to proliferative DR (Grade~4), is widely adopted in clinical
practice~\cite{wilkinson2003}. Although effective, manual grading is
time-intensive, costly, and subject to inter-grader variability---a limitation
particularly acute in low- and middle-income countries~\cite{krause2018}.

The emergence of deep learning (DL) has fundamentally transformed automated DR
grading. Landmark studies by Gulshan et al.~\cite{gulshan2016} and Ting et
al.~\cite{ting2017} demonstrated that CNNs could achieve diagnostic accuracy on par
with ophthalmologists. Subsequent research explored ResNet~\cite{he2016},
EfficientNet~\cite{tan2019}, and DenseNet~\cite{huang2017} architectures, achieving
substantial gains on benchmark challenges. A comprehensive survey by Zhu et
al.~\cite{zhu2024} covering 94 studies confirmed that transfer learning with ResNet
and VGGNet constitutes the dominant paradigm, yet two fundamental challenges remain
unresolved.

The \textit{first} is the scarcity of labeled training data. The APTOS~2019
benchmark, one of the largest publicly available graded fundus datasets, contains
only 3,662 labeled images---orders of magnitude smaller than datasets used to
pretrain general-purpose vision models. The \textit{second} challenge is the
ordinal nature of the DR grading scale. Standard cross-entropy loss treats all
misclassification errors equally, failing to encode the cost distinction between
confusing Grade~0 with Grade~1 versus Grade~0 with Grade~4. This misalignment with
the QWK evaluation metric represents a significant structural limitation, as
recently highlighted by Bappi et al.~\cite{bappi2025}.

In parallel, Vision Transformers (ViT)~\cite{dosovitskiy2021} have demonstrated
that self-attention mechanisms can model long-range spatial dependencies that
locally constrained convolutions may miss. Systematic comparisons confirm that
hybrid CNN-ViT models consistently outperform either architecture in isolation on DR
grading~\cite{zhang2025}. Knowledge distillation (KD)~\cite{hinton2015} offers a
compelling framework for the low-data challenge: a compact student network is
trained on soft probability distributions produced by a powerful teacher,
effectively inheriting richer learned representations. Islam et
al.~\cite{islam2023} validated multi-teacher KD for resource-constrained DR
classification, yet most existing approaches employ a single teacher and single
distillation temperature, and ordinal task structure is rarely integrated into the
distillation objective.

This paper proposes \textit{TS\_ConvNeXtTiny\_Residual}, addressing all four
concurrent limitations. The principal contributions are:
\begin{enumerate}
  \item A \textbf{hybrid teacher ensemble} combining a custom Baseline CNN-ViT
  hybrid and an Advanced ResNet50+ViT-B/16 model with spectral normalization via
  equal-weight prediction averaging, yielding a more robust soft-label target than
  any single teacher.
  \item A \textbf{multi-temperature distillation scheme} at $T = \{2.0, 4.0\}$
  enabling simultaneous transfer of fine-grained discriminative and coarse semantic
  knowledge across DR severity grades.
  \item An \textbf{ordinal CDF loss} that directly penalizes deviations from the
  correct cumulative grade distribution, aligning the training objective with the
  QWK metric.
  \item A \textbf{non-inferiority loss with scheduled margin} that enforces
  sample-wise performance parity between student and teacher ensemble throughout
  training, preventing early-training overfitting.
  \item Evaluation on APTOS~2019 demonstrating mean test QWK of $\mathbf{0.9186}$
  and accuracy of $\mathbf{85.47\%}$ across three random seeds, surpassing all
  baselines examined.
\end{enumerate}

% ── 2. Literature Review ──────────────────────────────────────────────────────
\section{Literature Review}

\subsection{Deep Learning for DR Grading}

The application of DL to automated DR grading has evolved rapidly since Gulshan et
al.~\cite{gulshan2016} and Ting et al.~\cite{ting2017} established that CNNs could
approach ophthalmologist-level accuracy. The comprehensive review by Zhu et
al.~\cite{zhu2024} confirms that transfer learning with ResNet and VGGNet
constitutes the dominant paradigm on APTOS~2019 and EyePACS benchmarks. On the
APTOS~2019 five-class task specifically, Farag et al.~\cite{farag2022} proposed a
DenseNet169 encoder augmented with CBAM attention, achieving 82\% accuracy and
QWK~=~0.888 with cross-entropy only. Shakibania et al.~\cite{shakibania2024}
introduced a dual-branch EfficientNet-B0+ResNet50 architecture demonstrating the
benefit of multi-backbone ensemble designs.

Recent work has increasingly focused on enriching feature representations. Tian et
al.~\cite{tian2023} proposed the FA+KC-Net combining fine-grained attention with
lesion knowledge extraction, achieving top-tier performance across APTOS~2019,
DDR, Messidor, and EyePACS. Romero-Ora\'a et al.~\cite{romero2024} developed an
attention-based framework separating attention on dark and bright retinal
structures. Despite these advances, the recurring limitation identified by Bappi et
al.~\cite{bappi2025} is the inadequate handling of the ordinal structure of DR
severity, motivating ordinal-aware training objectives of the kind proposed herein.

\subsection{CNN-ViT Hybrid Architectures}

The ViT, introduced by Dosovitskiy et al.~\cite{dosovitskiy2021}, demonstrated
that self-attention can model global spatial dependencies that locally constrained
convolutions may miss---capturing relationships between distant retinal regions such
as optic disc vascular patterns and macular lesions. Zhang et
al.~\cite{zhang2025} performed a systematic evaluation of CNN, ViT, and hybrid
architectures on DR classification: the best hybrid (CvT-13) achieved QWK~=~0.84
and AUC~=~0.93, and interpretability analysis via Grad-CAM and Attention-Rollout
confirmed that CNNs focus on fine-grained lesion textures whereas ViTs exhibit
broader contextual attention---directly motivating their combination. Ikram and
Imran~\cite{ikram2025} proposed ResViT FusionNet with LIME and Grad-CAM
explainability; Azzou et al.~\cite{azzou2025} found their ResNet50+ViT hybrid
outperformed both standalone models with 98\% average precision.

In this work, the structural diversity between the two teacher networks is
empirically validated in Section~\ref{sec:ensemble} (Table~\ref{tab:ensemble}):
Teacher~1 achieves test QWK~=~0.640 while Teacher~2 achieves QWK~=~0.901, and
their equal-weight ensemble reaches QWK~=~0.902, confirming that the two teachers
commit complementary errors whose combination yields a more reliable soft-label
signal.

\subsection{Knowledge Distillation for Medical Image Classification}

Knowledge distillation, first systematically formulated by Hinton et
al.~\cite{hinton2015}, trains a compact student by minimizing a composite loss
combining hard-label cross-entropy and KL-divergence between temperature-softened
teacher and student distributions. The key insight is that soft teacher
probabilities encode inter-class similarity---``dark knowledge''---absent from
one-hot labels. A comprehensive survey is provided by Gou et al.~\cite{gou2021}.

For DR grading, Luo et al.~\cite{luo2020} proposed a self-KD framework combining
self-distillation with CAM-attention, achieving AUC~=~0.965 and 92.9\% accuracy on
Messidor without an external teacher. Ju et al.~\cite{ju2021} proposed the SALL
framework leveraging KD jointly with multi-task learning across DR and AMD labels,
improving grading accuracy by 5.91\%. Islam et al.~\cite{islam2023} demonstrated
that multi-teacher KD---ResNet152V2+Swin Transformer into a lightweight Xception
student---can substantially outperform single-teacher approaches. Al~Shafi et
al.~\cite{alshafi2024} applied ViT-to-CNN distillation on APTOS achieving weighted
kappa~=~0.90.

A critical limitation of standard single-temperature KD is the inability to
simultaneously transfer both coarse semantic and fine-grained discriminative
knowledge. The present work addresses this through multi-temperature distillation at
$T = \{2.0, 4.0\}$, drawing on the insight established in the KD
literature~\cite{hinton2015,gou2021} that higher temperatures produce more uniform
soft distributions favorable for coarse knowledge transfer while lower temperatures
preserve sharper distributional peaks essential for distinguishing adjacent DR
severity grades.

The ordinal structure of DR grading further motivates specialized loss formulations.
Cao et al.~\cite{cao2020} proposed CORAL, reformulating ordinal regression as a set
of binary classification sub-problems. Liu et al.~\cite{liu2020} introduced
unimodal regularization that concentrates probability mass near the true ordinal
class. In this work, we adopt an ordinal CDF loss that directly penalizes
deviations between predicted and target cumulative distribution functions, providing
a differentiable ordinal regularizer explicitly aligned with the QWK metric.

\subsection{Teacher Ensemble and Multi-Teacher Distillation}

Ensemble learning improves predictive accuracy by aggregating diverse model
predictions. Yilmaz and Aiyengar~\cite{yilmaz2025} showed that cross-architecture
KD from a ViT teacher to a CNN student achieves 89\% classification accuracy on
fundus images with 97.4\% fewer parameters, retaining $\approx$93\% of the
teacher's diagnostic performance. Islam et al.~\cite{islam2023} further demonstrated
that a ResNet152V2+Swin Transformer teacher ensemble substantially outperformed
either teacher in isolation for DR classification.

Most multi-teacher approaches employ either simple averaging or confidence-weighted
combination. In this work, we investigated an entropy-gated mechanism but found
empirically that in the low-data setting (mean teacher confidence $\approx$~0.88)
the gate consistently collapsed to near-zero student weighting (mean gate value
$\approx$~0.09), yielding results statistically indistinguishable from equal-weight
averaging. This is consistent with the KD literature finding that gating mechanisms
require sufficient inter-teacher prediction uncertainty to provide
benefit~\cite{gou2021}.

\textbf{Research Gap.} The foregoing review identifies four concurrent limitations
no existing DR grading system has addressed in combination: (i)~reliance on a
single teacher whose soft labels may be miscalibrated in low-data
conditions---evidenced by Teacher~1 alone yielding QWK~=~0.640 (Table~\ref{tab:ensemble}); (ii)~fixed single-temperature distillation that cannot
simultaneously preserve fine-grained and coarse knowledge~\cite{hinton2015,gou2021}; (iii)~structural misalignment between cross-entropy
training and the QWK evaluation metric~\cite{bappi2025}; and (iv)~absence of a
non-inferiority mechanism to prevent student regression during early training. The
proposed \textit{TS\_ConvNeXtTiny\_Residual} framework addresses all four
limitations through a unified multi-teacher ensemble distillation pipeline with
ordinal-aware loss formulation.

% ── 3. Proposed Model ─────────────────────────────────────────────────────────
\section{Proposed Model}

\subsection{Overview}
\label{sec:overview}

The proposed framework is a teacher--student knowledge distillation system designed
for five-class DR severity grading under low-data conditions. Figure~\ref{fig:architecture} illustrates the complete pipeline. Given a fundus image of
size $224 \times 224 \times 3$, two independently trained CNN-ViT hybrid teacher
networks produce class probability distributions combined via equal-weight ensemble
averaging to yield a composite teacher prediction $\mathbf{p}_{teacher}$.
Concurrently, a ConvNeXt~Tiny student network~\cite{liu2022} processes the same
input to produce student logits $\mathbf{z}_s$. During training, the student is
optimized using a four-component loss integrating weighted cross-entropy,
multi-temperature distillation, ordinal CDF regularization, and a non-inferiority
constraint.

\begin{figure}[h]
  \centering
  \fbox{\parbox{0.9\columnwidth}{\centering\vspace{1.5cm}
  [Figure 1: Architecture diagram to be inserted]\vspace{1.5cm}}}
  \caption{The proposed TS\_ConvNeXtTiny\_Residual pipeline. Two CNN-ViT hybrid
  teachers produce an equal-weight ensemble soft-label target. The ConvNeXt~Tiny
  student is trained with a four-component loss and combined with the teacher
  ensemble at inference.}
  \label{fig:architecture}
\end{figure}

\subsection{Teacher Network 1: Baseline CNN-ViT Hybrid}

The first teacher is a dual-branch architecture combining CNN and ViT
representations through a confidence-gated fusion mechanism.

\textbf{CNN Branch.} The CNN branch processes the input through a stem:
Conv2d($3{\to}64$, $k{=}7$, $s{=}2$, $p{=}3$) $\to$ BatchNorm2d $\to$ ReLU $\to$
MaxPool2d($k{=}3$, $s{=}2$, $p{=}1$), followed by three residual-style stages
increasing channel dimensionality ($64{\to}128{\to}256$). Global average pooling
and a fully-connected layer ($256{\to}512$) yield $\mathbf{f}_{cnn} \in
\mathbb{R}^{512}$.

\textbf{ViT Branch.} A $16{\times}16$ patch embedding ($3{\to}256$) produces
$N{=}196$ tokens. Six transformer blocks with 8-head multi-head self-attention and
MLP ($256{\to}512{\to}256$) encode the token sequence. Mean pooling and linear
projection ($256{\to}512$) yield $\mathbf{f}_{vit} \in \mathbb{R}^{512}$.

\textbf{Confidence-Gated Fusion.} Scalar confidence scores $s_c = \sigma(W_c
\mathbf{f}_{cnn})$ and $s_t = \sigma(W_t \mathbf{f}_{vit})$ gate a two-layer MLP:
\begin{equation}
  \alpha = \sigma\!\left(\text{MLP}([s_c,\, s_t])\right), \quad
  \mathbf{f}_{fused} = \alpha\, \mathbf{f}_{cnn} + (1{-}\alpha)\, \mathbf{f}_{vit}
\end{equation}
A linear head ($512{\to}5$) produces logits $\mathbf{z}_{base}$.

\subsection{Teacher Network 2: Advanced CNN-ViT with Spectral Normalization}

A pretrained ResNet50 backbone and a pretrained ViT-B/16 encoder process the input;
their outputs ($2048$ and $768$ dimensions respectively) are concatenated to yield a
$2816$-dimensional joint feature. A spectral-normalized linear head~\cite{miyato2018}
maps this to five logits $\mathbf{z}_{adv}$. Spectral normalization constrains the
Lipschitz constant of the output layer, smoothing the decision boundary and improving
calibration---a property critical for the near-boundary confusion cases prevalent in
DR grading.

\subsection{Teacher Ensemble via Equal-Weight Fusion}
\label{sec:ensemble}

The composite teacher prediction is:
\begin{equation}
  \mathbf{p}_{teacher} = \frac{1}{2}\!\left[\text{softmax}(\mathbf{z}_{base}) +
  \text{softmax}(\mathbf{z}_{adv})\right]
\end{equation}

The teacher logit for distillation is $\mathbf{z}_t = \log(\mathbf{p}_{teacher})$.
Equal-weight fusion was selected over entropy-gated weighting based on empirical
findings in Section~\ref{sec:teacher_selection}: the gating mechanism provided no
measurable improvement in the high-confidence teacher regime (mean teacher
confidence $\approx$~0.88, mean gate value $\approx$~0.09), making the simpler
equal-weight formulation the preferred production strategy.

\subsection{Student Network: ConvNeXt Tiny}

A key design decision is the choice of student backbone capacity. Lightweight
CNN-ViT architectures (524K parameters, explored in the ablation study in
Section~\ref{sec:ablation}) are parameter-efficient but empirically insufficient
to fully exploit the multi-teacher soft labels on the five-class ordinal task.
Accordingly, we select ConvNeXt~Tiny~\cite{liu2022} as the student backbone: a
purely convolutional architecture incorporating depthwise convolutions, inverted
bottleneck blocks, and GELU activations. ConvNeXt~Tiny achieves competitive
ImageNet accuracy with substantially fewer FLOPs than ViT-S/16, making it
well-suited for resource-constrained clinical deployment~\cite{islam2023,yilmaz2025}.
The backbone is initialized with ImageNet-pretrained weights; the classification
head is replaced with a linear layer ($768{\to}5$) to produce student logits
$\mathbf{z}_s$.

\subsection{Multi-Component Training Loss}

The student training objective addresses four distinct failure modes of standard
distillation applied to ordinal medical grading: (1)~class imbalance biasing
predictions toward majority grades, (2)~single-temperature KL-divergence failing to
transfer both fine-grained and coarse knowledge simultaneously, (3)~training--evaluation
misalignment between cross-entropy and QWK, and (4)~student regression below teacher
performance during early training. The composite objective is:
\begin{equation}
  \mathcal{L} = \mathcal{L}_{main} + \lambda_{ni}\,\mathcal{L}_{ni} +
  \lambda_{distill}\,\mathcal{L}_{distill} + \lambda_{ord}\,\mathcal{L}_{ord}
\end{equation}

\textbf{Weighted Cross-Entropy.}
\begin{equation}
  \mathcal{L}_{main} = -\sum_{c=0}^{4} w_c\, y_c\, \log\hat{p}_c, \quad
  w_c = \frac{N}{C \cdot n_c}
\end{equation}

\textbf{Non-Inferiority Loss.}
\begin{equation}
  \mathcal{L}_{ni} = \frac{1}{N}\sum_{i=1}^{N}
  \max\!\left(0,\; \ell_s^{(i)} - \ell_t^{(i)} + m(\text{epoch})\right)
\end{equation}
where $m(\text{epoch})$ decreases linearly from 0.010 to 0.003 over 70 epochs and
$\lambda_{ni}: 0.80 \to 0.25$.

\textbf{Multi-Temperature Distillation Loss.}
\begin{equation}
  \mathcal{L}_{distill} = \sum_{k \in \{1,2\}} w_k\, T_k^2 \cdot
  \text{KL}\!\left(\text{softmax}\!\left(\tfrac{\mathbf{z}_s}{T_k}\right) \,\Big\|\,
  \text{softmax}\!\left(\tfrac{\mathbf{z}_t}{T_k}\right)\right)
\end{equation}
with $T_1{=}2.0$ ($w_1{=}0.6$) and $T_2{=}4.0$ ($w_2{=}0.4$). The $T_k^2$
scaling ensures gradient magnitude consistency across temperatures~\cite{hinton2015}.
The overall distillation weight is $\lambda_{distill} = 0.30$.

\textbf{Ordinal CDF Loss.}
\begin{equation}
  \mathcal{L}_{ord} = \frac{1}{4}\sum_{k=0}^{3} \left(\hat{F}_k - F_k^*\right)^2
\end{equation}
where $\hat{F}_k = \sum_{c=0}^{k} \hat{p}_c$ and $F_k^* = \mathbf{1}[y \leq k]$.
The ordinal weight is $\lambda_{ord} = 0.12$.

\subsection{Inference-Time Ensemble}

At inference, the final class probability is:
\begin{equation}
  \mathbf{p}_{final} = \frac{1}{2}\!\left[\mathbf{p}_{teacher} +
  \text{softmax}(\mathbf{z}_s)\right], \quad
  \hat{y} = \arg\max_c\; [\mathbf{p}_{final}]_c
\end{equation}

\subsection{Training Configuration}

All components are implemented in PyTorch on NVIDIA GPU hardware. The APTOS~2019
training set is partitioned using stratified random splitting into training (80\%,
$n{=}2{,}929$), validation (10\%, $n{=}366$), and test (10\%, $n{=}367$) subsets.
All final-stage experiments are repeated across three random seeds (42, 52, 62) and
reported as mean $\pm$ standard deviation. Hyperparameters are listed in
Table~\ref{tab:hyperparams}.

\begin{table}[h]
\centering
\caption{Training Hyperparameters}
\label{tab:hyperparams}
\small
\begin{tabular}{ll}
\toprule
\textbf{Hyperparameter} & \textbf{Value} \\
\midrule
Optimizer & AdamW \\
Learning rate & $2 \times 10^{-4}$ \\
Weight decay & $1 \times 10^{-4}$ \\
Batch size & 16 \\
Maximum epochs & 70 \\
LR scheduler & CosineAnnealingWarmRestarts ($T_0{=}10$, $T_{mult}{=}2$) \\
Early stopping patience & 10 epochs (val QWK) \\
Gradient clipping (max norm) & 1.0 \\
$\lambda_{distill}$ & 0.30 \\
$\lambda_{ordinal}$ & 0.12 \\
$\lambda_{ni}$ (scheduled) & $0.80 \to 0.25$ (linear) \\
NI margin (scheduled) & $0.010 \to 0.003$ (linear) \\
Distillation temperatures & $T_1{=}2.0$ ($w{=}0.6$), $T_2{=}4.0$ ($w{=}0.4$) \\
\bottomrule
\end{tabular}
\end{table}

% ── 4. Experimental Results ───────────────────────────────────────────────────
\section{Experimental Results}

\subsection{Dataset and Evaluation Protocol}

All experiments are conducted on the APTOS~2019 Blindness Detection
dataset~\cite{aptos2019}, which comprises 3,662 fundus photographs annotated on the
ICDR five-class scale: No DR (Grade~0, $n{=}1{,}805$), Mild ($n{=}370$), Moderate
($n{=}999$), Severe ($n{=}193$), and Proliferative DR ($n{=}295$). The primary
evaluation metric is QWK; secondary metrics include accuracy, macro-averaged
F1-score, and ROC-AUC (one-vs-rest).

\subsection{Teacher Architecture Selection}
\label{sec:teacher_selection}

\textbf{Stage~1: Standalone Backbone Baselines.} Nine backbone configurations were
evaluated under identical conditions. Results are reported in
Table~\ref{tab:backbones}. The custom Fusion CNN-ViT achieves the highest QWK
(0.923), confirming the well-established advantage of hybrid architectures for DR
grading~\cite{zhang2025}. ConvNeXt~Tiny achieves QWK~=~0.906 in standalone mode,
validating its selection as the student backbone.

\begin{table}[h]
\centering
\caption{Standalone Backbone Baselines on APTOS~2019 Test Set}
\label{tab:backbones}
\small
\begin{tabular}{lcccc}
\toprule
\textbf{Model} & \textbf{Acc (\%)} & \textbf{F1-Mac} & \textbf{QWK} & \textbf{AUC} \\
\midrule
Fusion CNN-ViT (custom) & 83.11 & 0.659 & \textbf{0.923} & 0.929 \\
Fusion ResNet+EfficientNet & 83.65 & 0.674 & 0.909 & 0.932 \\
EfficientNet-B0 & 82.56 & 0.658 & 0.908 & 0.897 \\
Fusion DenseNet+EfficientNet & 84.47 & 0.666 & 0.906 & 0.944 \\
ConvNeXt Tiny (standalone) & 79.02 & 0.648 & 0.906 & 0.947 \\
ResNet50 & 81.74 & 0.654 & 0.904 & 0.919 \\
DenseNet121 & 82.56 & 0.688 & 0.897 & 0.931 \\
ViT-B/16 & 82.56 & 0.631 & 0.896 & 0.928 \\
EfficientNet-B3 & 73.84 & 0.605 & 0.886 & 0.939 \\
\bottomrule
\end{tabular}
\end{table}

\textbf{Stage~2: Advanced Teacher Candidate Selection.} Four hybrid variants were
evaluated as candidates for Teacher~2. SpectralNorm was selected over the
marginally higher-QWK FiLM Fusion due to its superior calibration
properties~\cite{miyato2018}. Results are in Table~\ref{tab:advanced}.

\begin{table}[h]
\centering
\caption{Advanced Teacher Candidates on APTOS~2019 Test Set}
\label{tab:advanced}
\small
\begin{tabular}{lcccc}
\toprule
\textbf{Model} & \textbf{Acc (\%)} & \textbf{F1-W} & \textbf{QWK} & \textbf{Decision} \\
\midrule
FiLM Fusion & 84.20 & 0.828 & 0.901 & $\times$ \\
\textbf{SpectralNorm} & \textbf{83.11} & \textbf{0.825} & \textbf{0.901} & \checkmark~Teacher~2 \\
Deformable Token & 82.56 & 0.817 & 0.890 & $\times$ \\
Prototype Head & 83.38 & 0.818 & 0.886 & $\times$ \\
\bottomrule
\end{tabular}
\end{table}

\subsection{Teacher Ensemble Configuration}
\label{sec:teacher_ensemble}

Table~\ref{tab:ensemble} reports the four ensemble configurations evaluated.

\begin{table}[h]
\centering
\caption{Teacher Ensemble and Fusion Strategy Comparison}
\label{tab:ensemble}
\small
\begin{tabular}{lccccc}
\toprule
\textbf{Model} & \textbf{Val Acc} & \textbf{Val QWK} & \textbf{Test Acc} & \textbf{Test QWK} & \textbf{AUC} \\
\midrule
Teacher 1 (Baseline) & 64.85 & 0.744 & 63.22 & 0.640 & 0.904 \\
Teacher 2 (Advanced) & 82.29 & 0.882 & 83.11 & 0.901 & 0.964 \\
Equal-Weight (T1+T2) & 83.11 & 0.879 & 83.11 & 0.902 & 0.953 \\
Gated (T1+T2) & 81.47 & 0.878 & 84.47 & 0.907 & 0.956 \\
Student (standalone) & 5.18 & 0.000 & 5.18 & 0.000 & 0.588 \\
Equal-W Fusion (Ens+Student) & 83.11 & 0.880 & 83.11 & 0.897 & 0.953 \\
\bottomrule
\end{tabular}
\end{table}

The gate collapse analysis revealed a mean gate value $\bar{g} = 0.090$ with 86\%
of test samples yielding $g < 0.20$, confirming that equal-weight fusion is the
preferred strategy in high-confidence low-data regimes.

\subsection{Student Training Ablation}
\label{sec:ablation}

Table~\ref{tab:ablation} reports the four-step progressive ablation.

\begin{table*}[t]
\centering
\caption{Student Training Ablation --- APTOS~2019 Test Set (Progressive Development)}
\label{tab:ablation}
\begin{tabular}{llccccc}
\toprule
\textbf{Model} & \textbf{Key Addition} & \textbf{Stage} & \textbf{Acc (\%)} & \textbf{F1-Mac} & \textbf{QWK} & \textbf{Params} \\
\midrule
TS\_LiteCNNViT\_BaseResidual & CNN-ViT + ResidualDistill & screen & 85.01 & 0.701 & 0.909 & 524K \\
TS\_LiteCNNViT\_ConflictBrake & + ConflictBrake gate & screen & 85.29 & 0.702 & 0.912 & 524K \\
TS\_LiteCNNViT\_Conflict\_OrdinalDistill & + OrdinalDistill & final & $84.56 \pm 0.69$ & $0.687 \pm 0.017$ & $0.9155 \pm 0.0020$ & 524K \\
\textbf{TS\_ConvNeXtTiny\_Residual (proposed)} & \textbf{ConvNeXt Tiny backbone} & \textbf{final} & $\mathbf{85.47 \pm 0.31}$ & $\mathbf{0.703 \pm 0.009}$ & $\mathbf{0.9186 \pm 0.0010}$ & \textbf{27.8M} \\
\bottomrule
\end{tabular}
\end{table*}

\subsection{Comparison with State-of-the-Art}

Table~\ref{tab:sota} compares the proposed framework against published methods on
the APTOS~2019 five-class DR severity grading task.

\begin{table}[h]
\centering
\caption{Comparison with State-of-the-Art on APTOS~2019 (Five-Class Grading)}
\label{tab:sota}
\small
\begin{tabular}{lcccc}
\toprule
\textbf{Method} & \textbf{Acc (\%)} & \textbf{QWK} & \textbf{F1} & \textbf{Notes} \\
\midrule
Farag et al.~\cite{farag2022} & 82.0 & 0.888 & --- & CE only \\
Nazih et al.~\cite{nazih2023} & --- & --- & 0.825 & AUC=0.956 \\
Zhang et al.~\cite{zhang2025} & --- & 0.840 & --- & AUC=0.93 \\
Al~Shafi et al.~\cite{alshafi2024} & --- & 0.900\textsuperscript{\dag} & --- & SHAP \\
Ikram \& Imran~\cite{ikram2025} & 93.01 & 0.894\textsuperscript{\ddag} & --- & diff. split \\
Shakibania et al.~\cite{shakibania2024} & 89.6 & 0.930 & --- & augmented \\
Islam et al.~\cite{islam2023} & 99.0\textsuperscript{\P} & --- & --- & diff. split \\
\midrule
\textbf{Proposed (ours)} & $\mathbf{85.47 \pm 0.31}$ & $\mathbf{0.9186 \pm 0.0010}$ & \textbf{0.703} & no aug. \\
\bottomrule
\end{tabular}
\vspace{2pt}
{\small \textsuperscript{\dag}Weighted kappa. \textsuperscript{\ddag}Cohen's kappa, different APTOS split. \textsuperscript{\P}Different train--test split.}
\end{table}

\subsection{Computational Efficiency}

\begin{table}[h]
\centering
\caption{Computational Efficiency Comparison}
\label{tab:efficiency}
\small
\begin{tabular}{lccc}
\toprule
\textbf{Model} & \textbf{Train. Params} & \textbf{Total Params} & \textbf{Latency (ms)} \\
\midrule
Teacher 1 & $\sim$8.2M & $\sim$8.2M & --- \\
Teacher 2 & $\sim$113.8M & $\sim$113.8M & --- \\
TS\_LiteCNNViT & 524K & 114.7M & $5.41 \pm 0.01$ \\
\textbf{Proposed} & \textbf{27.8M} & \textbf{142.0M} & $\mathbf{6.70 \pm 0.01}$ \\
\bottomrule
\end{tabular}
\end{table}

% ── 5. Conclusion ─────────────────────────────────────────────────────────────
\section{Conclusion}

This paper presented TS\_ConvNeXtTiny\_Residual, a teacher--student knowledge
distillation framework for five-class DR severity grading under low-data conditions,
addressing the scarcity of expert-annotated training data and the failure of
standard cross-entropy losses to encode ordinal structure. The principal
contributions and their empirical validation are: (1)~a \textbf{hybrid teacher
ensemble} outperforming either teacher in isolation by up to 0.018 QWK; (2)~a
\textbf{multi-temperature distillation scheme} at $T \in \{2.0, 4.0\}$ enabling
simultaneous transfer of fine-grained and coarse knowledge; (3)~an \textbf{ordinal
CDF loss} yielding a consistent gain of $\approx$0.003 QWK over configurations
without ordinal regularization; (4)~a \textbf{non-inferiority loss with scheduled
margin} stabilizing learning in the low-data regime; and (5)~an empirical
demonstration of \textbf{gate collapse} under high-confidence teacher conditions,
providing a principled justification for equal-weight fusion. Evaluated across three
random seeds on APTOS~2019, the proposed system achieves mean test QWK of
$\mathbf{0.9186 \pm 0.0010}$ and accuracy of $\mathbf{85.47\% \pm 0.31\%}$,
surpassing all baselines examined.

\textbf{Limitations and Future Work.} The present evaluation is restricted to
APTOS~2019; multi-dataset evaluation across EyePACS, DDR, and Messidor would
strengthen generalizability claims. The deliberate exclusion of student-side
augmentation---methodologically justified for isolating distillation effects---represents quantifiable performance headroom. Incorporating asymmetric ordinal
penalties aligned with clinical triage priorities (under-grading proliferative DR
is clinically more severe than over-referring mild DR) is a promising direction.

% ── References ────────────────────────────────────────────────────────────────
\bibliographystyle{IEEEtran}

\begin{thebibliography}{99}

\bibitem{yau2012}
J.~W. Yau, S.~L. Rogers, R.~Kawasaki \textit{et al.}, ``Global prevalence and major risk factors of diabetic retinopathy,'' \textit{Diabetes Care}, vol.~35, no.~3, pp.~556--564, 2012. \url{https://doi.org/10.2337/dc11-1909}

\bibitem{idf2021}
International Diabetes Federation, \textit{IDF Diabetes Atlas}, 10th~ed. Brussels: IDF, 2021. \url{https://www.diabetesatlas.org/}

\bibitem{krause2018}
J.~Krause, V.~Gulshan, E.~Rahimy \textit{et al.}, ``Grader variability and the importance of reference standards for evaluating machine learning models for diabetic retinopathy,'' \textit{Ophthalmology}, vol.~125, no.~8, pp.~1264--1272, 2018.

\bibitem{gulshan2016}
V.~Gulshan, L.~Peng, M.~Coram \textit{et al.}, ``Development and validation of a deep learning algorithm for detection of diabetic retinopathy in retinal fundus photographs,'' \textit{JAMA}, vol.~316, no.~22, pp.~2402--2410, 2016.

\bibitem{ting2017}
D.~S.~W. Ting, C.~Y. Cheung, G.~Lim \textit{et al.}, ``Development and validation of a deep learning system for diabetic retinopathy and related eye diseases,'' \textit{JAMA}, vol.~318, no.~22, pp.~2211--2223, 2017.

\bibitem{he2016}
K.~He, X.~Zhang, S.~Ren, and J.~Sun, ``Deep residual learning for image recognition,'' in \textit{Proc. IEEE/CVF CVPR}, 2016, pp.~770--778.

\bibitem{tan2019}
M.~Tan and Q.~V. Le, ``EfficientNet: Rethinking model scaling for convolutional neural networks,'' in \textit{Proc. ICML}, 2019, pp.~6105--6114.

\bibitem{huang2017}
G.~Huang, Z.~Liu, L.~van~der~Maaten, and K.~Q. Weinberger, ``Densely connected convolutional networks,'' in \textit{Proc. IEEE/CVF CVPR}, 2017, pp.~4700--4708.

\bibitem{zhu2024}
S.~Zhu, C.~Xiong, Q.~Q. Zhong, and Y.~Yao, ``Diabetic retinopathy classification with deep learning via fundus images: A short survey,'' \textit{IEEE Access}, 2024. \url{https://doi.org/10.1109/access.2024.3361944}

\bibitem{bappi2025}
M.~I. Bappi, J.~A. Juthy, and K.~Kim, ``Deep learning-based diabetic retinopathy recognition and grading: Challenges, gaps, and an improved approach --- A survey,'' \textit{ICT Express}, 2025. \url{https://doi.org/10.1016/j.icte.2025.08.001}

\bibitem{dosovitskiy2021}
A.~Dosovitskiy, L.~Beyer, A.~Kolesnikov \textit{et al.}, ``An image is worth 16$\times$16 words: Transformers for image recognition at scale,'' in \textit{ICLR}, 2021. \url{https://arxiv.org/abs/2010.11929}

\bibitem{zhang2025}
W.~Zhang, V.~Belcheva, and T.~Ermakova, ``Interpretable deep learning for diabetic retinopathy: A comparative study of CNN, ViT, and hybrid architectures,'' \textit{Computers}, 2025. \url{https://doi.org/10.3390/computers14050187}

\bibitem{hinton2015}
G.~Hinton, O.~Vinyals, and J.~Dean, ``Distilling the knowledge in a neural network,'' \textit{arXiv preprint arXiv:1503.02531}, 2015.

\bibitem{islam2023}
N.~Islam, M.~M.~H. Jony, E.~Hasan \textit{et al.}, ``Toward lightweight diabetic retinopathy classification: A knowledge distillation approach for resource-constrained settings,'' \textit{Applied Sciences}, 2023. \url{https://doi.org/10.3390/app132212397}

\bibitem{liu2022}
Z.~Liu, H.~Mao, C.~Y. Wu \textit{et al.}, ``A ConvNet for the 2020s,'' in \textit{Proc. IEEE/CVF CVPR}, 2022, pp.~11976--11986.

\bibitem{farag2022}
M.~M. Farag, M.~A. Fouad, and A.~T. Abdel-Hamid, ``Automatic severity classification of diabetic retinopathy based on DenseNet and convolutional block attention module,'' \textit{IEEE Access}, 2022.

\bibitem{shakibania2024}
H.~Shakibania, S.~Raoufi, B.~Pourafkham \textit{et al.}, ``Dual branch deep learning network for detection and stage grading of diabetic retinopathy,'' \textit{Biomed. Signal Process. Control}, 2024.

\bibitem{tian2023}
M.~Tian, H.~Wang, Y.~Sun \textit{et al.}, ``Fine-grained attention \& knowledge-based collaborative network for diabetic retinopathy grading,'' \textit{Heliyon}, 2023.

\bibitem{romero2024}
R.~Romero-Ora\'a, M.~Herrero-Tudela, M.~I. L\'opez \textit{et al.}, ``Attention-based deep learning framework for automatic fundus image processing,'' \textit{Comput. Methods Programs Biomed.}, 2024.

\bibitem{aptos2019}
APTOS 2019 Blindness Detection, Kaggle Competition Dataset, 2019. \url{https://www.kaggle.com/c/aptos2019-blindness-detection}

\bibitem{ikram2025}
A.~Ikram and A.~Imran, ``ResViT FusionNet model: An explainable AI-driven approach for automated grading of diabetic retinopathy,'' \textit{Comput. Biol. Med.}, 2025.

\bibitem{azzou2025}
S.~Azzou, D.~Boukredera, and S.~Baouz, ``A hybrid CNN-Transformer approach for precise three-class diabetic retinopathy classification,'' \textit{Jordanian J. Comput. Inf. Technol.}, 2025.

\bibitem{niranjan2025}
K.~Niranjan \textit{et al.}, ``A hybrid vision transformer and graph neural network model for diabetic retinopathy detection,'' \textit{J. Artif. Intell. Technol.}, 2025.

\bibitem{nazih2023}
W.~Nazih, A.~O. Aseeri, O.~Atallah, and S.~El-Sappagh, ``Vision transformer model for predicting the severity of diabetic retinopathy,'' \textit{IEEE Access}, 2023.

\bibitem{goh2024}
J.~H.~L. Goh, E.~Ang, S.~Srinivasan \textit{et al.}, ``Comparative analysis of vision transformers and CNNs in detecting referable diabetic retinopathy,'' \textit{Ophthalmol. Sci.}, 2024.

\bibitem{gou2021}
J.~Gou, B.~Yu, S.~J. Maybank, and D.~Tao, ``Knowledge distillation: A survey,'' \textit{Int. J. Comput. Vis.}, vol.~129, pp.~1789--1819, 2021.

\bibitem{luo2020}
L.~Luo, D.~Xue, and X.~Feng, ``Automatic diabetic retinopathy grading via self-knowledge distillation,'' \textit{Electronics}, 2020.

\bibitem{ju2021}
L.~Ju, X.~Wang, X.~Zhao \textit{et al.}, ``Synergic adversarial label learning for grading retinal diseases via knowledge distillation,'' \textit{IEEE J. Biomed. Health Inform.}, 2021.

\bibitem{moya2024}
E.~Moya-Albor, A.~Lopez-Figueroa, S.~Jacome-Herrera \textit{et al.}, ``Computer-aided diagnosis of diabetic retinopathy lesions based on knowledge distillation,'' \textit{Mathematics}, 2024.

\bibitem{wang2024}
Y.~Wang, Q.~Hou, P.~Cao \textit{et al.}, ``Lesion-aware knowledge distillation for diabetic retinopathy lesion segmentation,'' \textit{Appl. Intell.}, 2024.

\bibitem{alshafi2024}
A.~Al~Shafi, M.~H. Shawon, N.~Shahid \textit{et al.}, ``Enhancing diabetic retinopathy detection through transformer based knowledge distillation,'' in \textit{Proc. IJCNN}, 2024.

\bibitem{cao2020}
W.~Cao, V.~Mirjalili, and S.~Raschka, ``Rank consistent ordinal regression for neural networks,'' \textit{Pattern Recognit. Lett.}, vol.~140, pp.~325--331, 2020.

\bibitem{liu2020}
X.~Liu, F.~Fan, L.~Kong \textit{et al.}, ``Unimodal regularized neuron stick-breaking for ordinal classification,'' \textit{Neurocomputing}, vol.~388, pp.~34--44, 2020.

\bibitem{yilmaz2025}
B.~Yilmaz and A.~Aiyengar, ``Cross-architecture knowledge distillation for retinal fundus image anomaly detection,'' \textit{arXiv preprint arXiv:2506.18220}, 2025.

\bibitem{miyato2018}
T.~Miyato, T.~Kataoka, M.~Koyama, and Y.~Yoshida, ``Spectral normalization for generative adversarial networks,'' in \textit{ICLR}, 2018.

\bibitem{wilkinson2003}
C.~P. Wilkinson, F.~L. Ferris~III, R.~E. Klein \textit{et al.}, ``Proposed international clinical diabetic retinopathy and diabetic macular edema disease severity scales,'' \textit{Ophthalmology}, vol.~110, no.~9, pp.~1677--1682, 2003.

\end{thebibliography}

\end{document}
